\chapter*{Resumen}
\addcontentsline{toc}{chapter}{Resumen}

DePaso es una plataforma digital de logística urbana colaborativa orientada al Área Metropolitana de Buenos Aires (AMBA), que conecta remitentes ---PyMEs, comercios y particulares--- con transportistas dispuestos a llevar encomiendas en trayectos que ya realizan habitualmente. El proyecto introduce un modelo híbrido que combina transporte dedicado con transporte colaborativo, aprovechando la capacidad ociosa de vehículos en circulación para reducir costos, emisiones de CO\textsubscript{2} y tiempos de entrega.

La plataforma contempla cuatro componentes técnicos principales: (1)~un clasificador de carga basado en visión computacional que determina el tipo y tamaño del paquete a partir de una fotografía; (2)~un algoritmo de asignación multivariable que pondera compatibilidad geoespacial, desvío, tipo de carga, reputación del transportista y ventana horaria; (3)~seguimiento del envío en tiempo real; y (4)~cálculo del ahorro de CO\textsubscript{2} respecto al escenario de transporte dedicado, basado en factores de emisión del IPCC.

{\color{nuevo}Este documento corresponde al avance del 50\% del Proyecto Final de Ingeniería y presenta el marco teórico que sustenta las decisiones de diseño, el estado del arte y el análisis competitivo, los resultados de la investigación de usuarios (una encuesta con 145 respuestas de residentes y trabajadores del AMBA, complementada con entrevistas semiestructuradas), el análisis de requerimientos con sus casos de uso, y la metodología de desarrollo: arquitectura, tecnologías, modelo de datos y validación del sistema. Los hallazgos cuantitativos evidencian alta disposición al uso de envíos colaborativos (69,8\%), fuerte preferencia por la transparencia en el precio (el 75\% señala el costo como criterio principal) y un gap de precios accionable entre la disposición a pagar de los remitentes (\$3.000--\$6.000 por un paquete chico) y la compensación mínima aceptable de los transportistas (\$2.500--\$5.000), lo que valida la viabilidad económica del modelo propuesto.}

\textbf{Palabras clave:} logística colaborativa, crowdsourced delivery, última milla, mercados de dos lados, visión computacional, AMBA.
